\documentclass[11pt]{article}
\usepackage[margin=0.85in]{geometry}
\usepackage{amsmath,amssymb,amsthm,booktabs,hyperref}
\hypersetup{colorlinks=true,urlcolor=blue,linkcolor=blue}
\newtheorem{theorem}{Theorem}
\newtheorem{lemma}{Lemma}
\title{The full Burnside--Stirling--Gautschi bound\\
\large MathDB 372858: solution of Schmidlin's Conjecture 37}
\author{Principia Math In Progress}
\date{12 September 2026}
\begin{document}
\maketitle

\section*{Task 1. Prove the bound on its full stated domain}
For \(x\geq0\) and \(0\leq\alpha\leq1\), put
\[
 \widehat\Pi(x,\alpha)=\Gamma(x+1)(x+1)^\alpha,\qquad
 S_{1/2}(y)=\sqrt{2\pi}(y+\tfrac12)^{y+1/2}e^{-y-1/2},
\]
and \(\widehat m_{1/2}(x,\alpha)=
\log\bigl(\widehat\Pi(x,\alpha)/S_{1/2}(x+\alpha)\bigr)\).
These are the definitions of \cite{Schmidlin}, pp.~3--4 and 10.
The domain \(\widehat D_{1/2}\) is exactly
\([0,\infty)\times[0,1]\). Conjecture 37, pp.~28--29, asks
to remove the restriction \(x\geq1\) in Corollary 32.

\begin{theorem}
For every \(x\geq0\) and \(\alpha\in[0,1]\),
\begin{equation}\label{target}
 \widehat m_{1/2}(x,\alpha)
 < \frac{1}{12(x+\alpha)+18-12e^{1/48}}.
\end{equation}
In particular, Conjecture 37 holds, with strict inequality.
\end{theorem}

We first record the only special-function estimate needed.
\begin{lemma}\label{binet}
For every \(u>0\), the Stirling remainder
\[
 r(u)=\log\Gamma(u)-(u-\tfrac12)\log u+u-\tfrac12\log(2\pi)
\]
satisfies
\begin{equation}\label{remainder}
 r(u)<B(u):=\frac1{12u}-\frac1{360u^3}+\frac1{1260u^5}.
\end{equation}
\end{lemma}
\begin{proof}
Binet's formula \cite[Eq.~5.9.10]{DLMF} gives
\[
 r(u)=2\int_0^\infty\frac{\arctan(t/u)}{e^{2\pi t}-1}\,dt.
\]
The exact identity
\[
 z-\frac{z^3}{3}+\frac{z^5}{5}-\arctan z
 =\int_0^z\frac{s^6}{1+s^2}\,ds>0\qquad(z>0)
\]
holds on the whole positive axis. Also,
\[
 \int_0^\infty\frac{t^j}{e^{2\pi t}-1}\,dt
 =\frac{j!\,\zeta(j+1)}{(2\pi)^{j+1}}
 =\begin{cases}1/24&j=1,\\1/240&j=3,\\1/504&j=5.\end{cases}
\]
The first equality follows by expanding the positive geometric series
and integrating termwise. Integrating the strict arctangent bound proves
\eqref{remainder}; every integral involved is finite.
\end{proof}

\newpage
\begin{proof}[Proof of the theorem]
Set \(u=x+1\geq1\), \(\delta=\alpha-\tfrac12\in[-\tfrac12,\tfrac12]\),
and \(d=12(e^{1/48}-1)\). The elementary inequalities
\(1+t<e^t<(1-t)^{-1}\), for \(0<t<1\), imply
\begin{equation}\label{d}
 \frac14<d<\frac{12}{47}<\frac13.
\end{equation}
Direct substitution in the definitions gives
\[
 \widehat m_{1/2}(x,\alpha)
 =r(u)+\delta-(u+\delta)\log(1+\delta/u).
\]
Thus we must show that
\[
 F_u(\delta):=r(u)+\delta-(u+\delta)\log(1+\delta/u)
             -\frac1{12(u+\delta)-d}<0.
\]
All denominators are positive, since \(u+\delta\geq\tfrac12\).
Differentiation yields
\begin{align}
 F'_u(\delta)&=-\log(1+\delta/u)+\frac{12}{[12(u+\delta)-d]^2},
                                                        \label{first}\\
 F''_u(\delta)&=-\frac1{u+\delta}
               -\frac{288}{[12(u+\delta)-d]^3}<0.         \label{second}
\end{align}
In particular \(F'_u(\delta)>0\) when \(\delta\leq0\). At
\(\delta=1/12\), use \(\log(1+t)>t/(1+t)\) for \(t>0\).
For \(u=1+y\), \(y\geq0\), one has
\[
 (12u+\tfrac23)^2-12(12u+1)
 =\frac{40}{9}+160y+144y^2>0.
\]
Together with \(d<1/3\), this shows
\[
 \frac{12}{(12u+1-d)^2}
 <\frac{12}{(12u+2/3)^2}
 <\frac1{12u+1}
 <\log(1+1/(12u)).
\]
Hence \(F'_u(1/12)<0\). Strict concavity now places the global
maximum of \(F_u\) on \([-\tfrac12,\tfrac12]\) inside \((0,1/12)\).

On \(0\leq\delta\leq1/12\), let
\[
 v=u+\tfrac1{12},\qquad
 K(u)=\frac1v+\frac1{6v^3}
     =\frac{12(144u^2+24u+25)}{(12u+1)^3}>0.
\]
Because \(d>0\) and \(u+\delta\leq v\), \eqref{second} implies
\(F''_u(\delta)\leq-K(u)\). Integrating twice and completing a square,
\[
 F_u(\delta)\leq F_u(0)+F'_u(0)\delta-\tfrac12K(u)\delta^2
 \leq F_u(0)+\frac{F'_u(0)^2}{2K(u)}.
\]
Lemma \ref{binet} and \eqref{d} give
\[
 F_u(0)<B(u)-\frac1{12u-1/4},\qquad
 0<F'_u(0)<\frac{12}{(12u-1/3)^2}.
\]
Consequently the global maximum of \(F_u\) is strictly less than
\(-R(u)\), where
\begin{equation}\label{rational}
 R(u)=\frac1{12u-1/4}-\frac1{12u}+\frac1{360u^3}
      -\frac1{1260u^5}-\frac{72}{(12u-1/3)^4K(u)}.
\end{equation}

\newpage
It remains to prove one rational inequality, \(R(u)>0\) for \(u\geq1\).
Write \(Q(u)=144u^2+24u+25\). Exact simplification gives
\[
 R(u)=\frac{N(u)}
 {2520u^5(36u-1)^4(48u-1)Q(u)},
\]
where the numerator can be formed without any division:
\begin{align*}
 N(u)={}&210u^4(36u-1)^4Q(u)
       +7u^2(36u-1)^4(48u-1)Q(u)\\
       &-2(36u-1)^4(48u-1)Q(u)
       -1224720u^5(48u-1)(12u+1)^3.
\end{align*}
For an explicit positivity certificate, write
\(N(1+y)=\sum_{j=0}^{10}c_jy^j\). Its integer coefficients are
\[
\begin{array}{r|r@{\qquad}r|r}
j&c_j&j&c_j\\ \hline
0&2417815645&6&8758379206944\\
1&130284291970&7&5353548256512\\
2&968629141445&8&2115822352896\\
3&3398691532320&9&490420998144\\
4&7091567419890&10&50791587840\\
5&9614662445664&&
\end{array}
\]
Every coefficient is strictly positive, so \(N(u)>0\) for \(u\geq1\).
The denominator is positive there as well. This proves \(R(u)>0\),
then \(F_u(\delta)<0\) on the full required interval, and finally
\eqref{target}.
\end{proof}

\noindent\textbf{Certificate and scope.}
The accompanying \texttt{verify\_polynomial.py} builds the numerator
from its four displayed summands using exact integers, checks the
translation \(u=1+y\) by both binomial expansion and Horner substitution,
and verifies the rational identity at exact rational points.
The displayed polynomial identity and positive coefficients certify
the infinite range; numerical sampling is not used as a proof.
The argument includes \(x=0\), both endpoints of \([0,1]\), and
arbitrarily large \(x\).

\section*{Tasks 2--5. Further directions after the full solve}
These are proposed tasks, not additional solved claims; each should
receive a fresh literature and public-claim check before being attacked.
\begin{description}
\item[Task 2.] Determine the largest constant \(C\) for which
\(\widehat m_{1/2}(x,\alpha)\leq[12(x+\alpha)+C]^{-1}\)
holds on the full domain, with a positive denominator, and identify
all extremizers or limiting extremizers.
\item[Task 3.] For the constant in \eqref{target}, analyze the unique
maximizing \(\delta\) of \(F_u\): determine its monotonicity in \(u\)
and derive an asymptotic expansion with an explicit remainder.
\item[Task 4.] Extend the same inequality, wherever its denominator is
positive, to as much of the natural Gamma-function domain
\(x>-1,\ 0\leq\alpha\leq1,\ x+\alpha>-\tfrac12\) as possible.
\item[Task 5.] Find a simpler positive expression below \(R(u)\) for
all \(u\geq1\), preserving its leading asymptotic term, to replace
the degree-ten coefficient certificate by a shorter quantitative bound.
\end{description}

\begin{thebibliography}{9}
\bibitem{Schmidlin}
M. Schmidlin, \emph{On Gautschi \& Stirling Identities, Asymptotics
and Inequalities for the Pi (or Gamma) Function},
arXiv:2601.21906v1 (2026),
\url{https://arxiv.org/abs/2601.21906}.
\bibitem{DLMF}
NIST Digital Library of Mathematical Functions,
Binet's formula, Eq.~5.9.10,
\url{https://dlmf.nist.gov/5.9#E10}.
\end{thebibliography}
\end{document}
